%===============================================================================
% DOCUMENT
%===============================================================================

%% Document class
\documentclass[a4paper,12pt]{scrreprt}

%% Include packages
\usepackage{packages}

\begin{document}

%% Include custom commands
\include{commands}

\pagenumbering{gobble}

%% Build cover
\definecolor{titlepagecolor}{RGB}{245,130,32}

%==========================================================================
% COLORED BAR ON THE LEFT SIDE
%==========================================================================

\backgroundsetup{
    scale=1,
    angle=0,
    opacity=1,
    contents={
        \begin{tikzpicture}[remember picture,overlay]
            \path [fill=titlepagecolor] (-10.5,-15) rectangle ++ (5,30);
            \node[color=white] at (-6.60,-11.5) {\bfseries {\fontsize{120}{60} \textsf{I}}};
            \node[color=titlepagecolor] at (-4.20,-11.5) {\bfseries {\fontsize{120}{60} \textsf{A}}};
        \end{tikzpicture}
    }
}

%==========================================================================
% COVER PAGE CONTENT
%==========================================================================

\title{\LARGE{Otimização de Distribuição Alimentar numa Catástrofe}}

\author{
    Flávia Alexandra da Silva Araújo (A96587)\\ \quad
    Joshua David Amaral Moreira (A105684)\\ \quad
    Miguel Fernandes Barbosa (A104451)\\ \quad
    Miguel Torres Carvalho (A95485)\\ \quad
}

%% Date
\date{\today}

%% Course
\newcommand{\Course}{Licenciatura em Engenharia Informática}

%% Department
\newcommand{\Department}{Escola de Engenharia}

%% UniName
\newcommand{\UniName}{Universidade do Minho}

%% UniPic
\newcommand{\UniPic}{\includegraphics[width=120pt]{img/eeum.png}}

%% University
\newcommand{\University}{
    \begin{flushleft}
         \UniPic   
    \end{flushleft}
    \textcolor{gray}{\small\textbf{\textsf{\UniName}}}\par
    \textcolor{gray!80!white}{\small{\textsf{\Department}}}\par
    \textcolor{gray!70!white}{\small{\textsf{\Course}}}
}

%% UC
\newcommand{\UC}{
    \begin{flushleft}
        \par\textcolor{titlepagecolor}{\Large\textbf{\textsf{Inteligência Artificial}}}
    \end{flushleft}
}

%% Project Phase
\newcommand{\SubTitle}{
    \begin{flushleft}
        \large\textbf{Resolução de Problemas com Algoritmos de Procura}
    \end{flushleft}
}

%% Group Info
\newcommand{\GroupInfo}{\par Grupo 21}

%% GitHub Repo
\newcommand{\GitHubRepo}{\par\url{https://github.com/flaviaraujo/project-IA}}

%% School Year
\newcommand{\SchoolYear}{
    \par\small{\textsf{Ano Letivo de 2024/2025}}
}

%% Define new command to show title, author and date
\makeatletter
\let\Title\@title
\let\Author\@author
\let\Date\@date
\makeatother

%==========================================================================
% BEGIN COVER PAGE
%==========================================================================

%% Make cover page
\newcommand{\makecover}{

%% Removes page number on footer
\thispagestyle{empty}

%% No indentation
\setlength{\parindent}{0em}

%% Put Background defined on \backgroundsetup, in this page
\BgThispage

%% Changing geometry to prevent overlay with text
%% At the end of back cover, geometry is default with \restoregeometry
\newgeometry{top=4cm,left=6cm,right=2cm,bottom=2cm}

%% builds university info defined previously
\University
\vspace{1cm}
\UC
\SubTitle
\SchoolYear

\vspace*{4cm}
%% bigger space (i think its the default one) between paragraphs
\setlength{\parskip}{1em}

%% builds title info defined previously
\par\textbf{\textsf{\huge\Title}}
\vspace{1cm}
%% builds author(s) info defined previously
\par\Author

\vspace{0.5cm}

%% builds date info defined previously
\par\Date
\restoregeometry
\pagebreak

}

%==========================================================================
% END COVER PAGE
%==========================================================================

\makecover

%% Default geometry
\newgeometry{top=3cm,left=3cm,right=3cm,bottom=4cm}

%% Save default geometry
\savegeometry{default}

%% Load default geometry with:
% \loadgeometry{default}


%===============================================================================
% BEGIN ABSTRACT PAGE
%===============================================================================

\renewenvironment{abstract}
 {\par\noindent\textbf{\Large\abstractname}\par\bigskip}
 {}

\begin{flushleft}
\begin{abstract}
    <<O resumo tem como objectivo descrever de forma sucinta o trabalho realizado. Deverá conter uma pequena introdução, seguida por uma breve descrição do trabalho realizado e terminando com uma indicação sumária do seu estado final.>>
    \par \textbf{Área de Aplicação}: <<Identificação da Área de trabalho.>>
    \par \textbf{Palavras-Chave}: <<Conjunto de palavras-chave que permitirão referenciar domínios de conhecimento, tecnologias, estratégias, etc., directa ou indirectamente referidos no relatório.>>
\end{abstract}
\end{flushleft}

\pagebreak

%===============================================================================
% END ABSTRACT PAGE
%===============================================================================


%===============================================================================
% BEGIN INDEXES PAGES
%===============================================================================

%% Changes table of content name
\renewcommand{\contentsname}{Índice}
\renewcommand{\listfigurename}{Índice de Figuras}
\renewcommand{\listtablename}{Índice de Tabelas}
\renewcommand{\lstlistlistingname}{Índice de \textit{Snippets}}

\tableofcontents
\pagebreak

\listoffigures
\pagebreak

\listoftables
\pagebreak

\lstlistoflistings
\pagebreak


%===============================================================================
% END INDEXES PAGES
%===============================================================================

\pagenumbering{arabic}

%===============================================================================
% BEGIN DESCRIÇÃO DO PROBLEMA
%===============================================================================

\chapter{Descrição do Problema}


%===============================================================================
% END DESCRIÇÃO DO PROBLEMA
%===============================================================================

%===============================================================================
% BEGIN FORMULAÇÃO DO PROBLEMA
%===============================================================================

\chapter{Formulação do Problema}

\begin{itemize}
    \item\textbf{Tipo:} Problema de Procura de Estado Único (?)
    \item\textbf{Estado Inicial:} Todos os centros de distribuição, veículos e suprimentos disponíveis são conhecidos.
     As localizações e necessidades das zonas afetadas são identificadas, com diferentes graus de urgência.
    \item\textbf{Teste Objetivo:} 
        \item\textbf{Operadores (Ações Disponíveis):} 
        \begin{itemize}
                \item\textit{Mover Veículo:} Deslocar um veículo de uma zona para outra, respeitando as limitações de combustível e condições da rota.
                \item\textit{Carregar/Descarregar Suprimentos:} Realizar operações de carga e descarga nos veículos, levando em consideração as capacidades.
                \item\textit{Reabastecer Combustível:} Realizar paragens no percurso para reabastecimento dos veículos. (?)
        \end{itemize}
    \item\textbf{Custo da Solução:} O custo é calculado com base na distância percorrida, tempo de viagem, uso de combustível e deterioração dos suprimentos perecíveis.
    \end{itemize}

%===============================================================================
% END FORMULAÇÃO DO PROBLEMA
%===============================================================================

%===============================================================================
% BEGIN TAREFAS REALIZADAS E DECISÕES TOMADAS
%===============================================================================

\chapter{Tarefas Realizadas e Decisões Tomadas}

%===============================================================================
% END TAREFAS REALIZADAS E DECISÕES TOMADAS
%===============================================================================

%===============================================================================
% BEGIN CONCLUSÕES E DISCUSSÃO DOS RESULTADOS OBTIDOS
%===============================================================================

\chapter{Conclusões e Discussão dos Resultados Obtidos}

%===============================================================================
% END CONCLUSÕES E DISCUSSÃO DOS RESULTADOS OBTIDOS
%===============================================================================



%===============================================================================
% BEGIN REFERÊNCIAS
%===============================================================================

%% Change bibliography name from “Bibliografia” to “Referências”
\renewcommand\bibname{Referências}

\begin{thebibliography}{9}

\bibitem{RussellNorvig}
Russell, S., \& Norvig, P. (2009). \textit{Artificial Intelligence - A Modern Approach} (3rd ed.). Pearson Education. ISBN-13: 9780136042594.

\bibitem{CostaSimoes}
Costa, E., \& Simões, A. (2008). \textit{Inteligência Artificial - Fundamentos e Aplicações}. FCA. ISBN: 978-972-722-34.

\end{thebibliography}

%% Add bibliografia to index
\addcontentsline{toc}{chapter}{Referências}

%===============================================================================
% END REFERÊNCIAS
%===============================================================================





\end{document}
